\begin{bt}%[9D6H18-2]
% nguon: Đề thi tham khảo Tuyển sinh 10 TP.HCM – Phụ lục Sở GD&ĐT TP.HCM – Bài 1
Cho hàm số $y = \dfrac{x^2}{2}$ có đồ thị $(P)$.
\begin{enumerate}[a)]
	\item Vẽ đồ thị $(P)$ của hàm số trên.
	\item Tìm những điểm $M$ thuộc $(P)$ có tung độ và hoành độ bằng nhau.
\end{enumerate}
\loigiai{
\begin{enumerate}[a)]
	\item Bảng giá trị tương ứng giữa $x$ và $y$ của hàm số $y = \dfrac{x^2}{2}$:
	\begin{center}
	\begin{tabular}{|c|c|c|c|c|c|}
	\hline
	$x$ & $-4$ & $-2$ & $0$ & $2$ & $4$ \\
	\hline
	$y = \dfrac{x^2}{2}$ & $8$ & $2$ & $0$ & $2$ & $8$ \\
	\hline
	\end{tabular}
	\end{center}
	Đồ thị $(P)$ là một đường parabol đi qua gốc toạ độ $O(0; 0)$, nhận trục tung $Oy$ làm trục đối xứng và nằm phía trên trục hoành.
	\item Điểm $M(x_M; y_M)$ thuộc $(P)$ có tung độ và hoành độ bằng nhau thoả mãn:
	$y_M = x_M \Leftrightarrow \dfrac{x_M^2}{2} = x_M \Leftrightarrow x_M^2 - 2x_M = 0 \Leftrightarrow x_M(x_M - 2) = 0 \Leftrightarrow \left[\begin{aligned} x_M &= 0 \Rightarrow y_M = 0 \\ x_M &= 2 \Rightarrow y_M = 2. \end{aligned}\right.$\\
	Vậy có hai điểm thoả mãn yêu cầu đề bài là $M_1(0; 0)$ và $M_2(2; 2)$.
\end{enumerate}
}
\end{bt}

\begin{bt}%[9D6H20-2]
% nguon: Đề thi tham khảo Tuyển sinh 10 TP.HCM – Phụ lục Sở GD&ĐT TP.HCM – Bài 2
Cho phương trình $2x^2 - 5x + 1 = 0$.
\begin{enumerate}[a)]
	\item Chứng minh phương trình trên có hai nghiệm phân biệt.
	\item Không giải phương trình, hãy tính giá trị của biểu thức $A = x_1(x_1 + 2024) + x_2(x_2 + 2025) - x_2$.
\end{enumerate}
\loigiai{
\begin{enumerate}[a)]
	\item Xét phương trình $2x^2 - 5x + 1 = 0$, ta có:
	$\Delta = b^2 - 4ac = (-5)^2 - 4\cdot 2\cdot 1 = 25 - 8 = 17 > 0$.\\
	Vì $\Delta > 0$ nên phương trình luôn có hai nghiệm phân biệt $x_1, x_2$.
	\item Theo định lí Viète, ta có:
	$\begin{cases} x_1 + x_2 = -\dfrac{b}{a} = \dfrac{5}{2} \\ x_1 x_2 = \dfrac{c}{a} = \dfrac{1}{2}. \end{cases}$\\
	Biến đổi biểu thức $A$:
	$A = x_1^2 + 2024x_1 + x_2^2 + 2025x_2 - x_2 = (x_1^2 + x_2^2) + 2024(x_1 + x_2) = (x_1 + x_2)^2 - 2x_1 x_2 + 2024(x_1 + x_2)$.\\
	Thay $x_1 + x_2 = \dfrac{5}{2}$ và $x_1 x_2 = \dfrac{1}{2}$ vào biểu thức $A$, ta được:
	$A = \left(\dfrac{5}{2}\right)^2 - 2\cdot\dfrac{1}{2} + 2024\cdot\dfrac{5}{2} = \dfrac{25}{4} - 1 + 5060 = \dfrac{21}{4} + 5060 = \dfrac{20261}{4} = 5065{,}25$.\\
	Vậy $A = \dfrac{20261}{4} = 5065{,}25$.
\end{enumerate}
}
\end{bt}

\begin{bt}%[9D8H26-3]
% nguon: Đề thi tham khảo Tuyển sinh 10 TP.HCM – Phụ lục Sở GD&ĐT TP.HCM – Bài 3
Biên độ nhiệt là khoảng cách chênh lệch giữa nhiệt độ cao nhất và nhiệt độ thấp nhất trong cùng một khoảng thời gian nhất định của cùng một vùng địa lí. Số liệu nhiệt độ ($^\circ\text{C}$) các ngày trong một tuần tại Thành phố Hồ Chí Minh được ghi nhận như sau:
\begin{itemize}
	\item Thứ Hai: cao nhất $36^\circ\text{C}$, thấp nhất $25^\circ\text{C}$ (biên độ $11^\circ\text{C}$).
	\item Thứ Ba: cao nhất $37^\circ\text{C}$, thấp nhất $26^\circ\text{C}$ (biên độ $11^\circ\text{C}$).
	\item Thứ Tư: cao nhất $35^\circ\text{C}$, thấp nhất $24^\circ\text{C}$ (biên độ $11^\circ\text{C}$).
	\item Thứ Năm: cao nhất $36^\circ\text{C}$, thấp nhất $25^\circ\text{C}$ (biên độ $11^\circ\text{C}$).
	\item Thứ Sáu: cao nhất $35^\circ\text{C}$, thấp nhất $22^\circ\text{C}$ (biên độ $13^\circ\text{C}$).
	\item Thứ Bảy: cao nhất $36^\circ\text{C}$, thấp nhất $23^\circ\text{C}$ (biên độ $13^\circ\text{C}$).
	\item Chủ Nhật: cao nhất $34^\circ\text{C}$, thấp nhất $27^\circ\text{C}$ (biên độ $7^\circ\text{C}$).
\end{itemize}
\begin{enumerate}[a)]
	\item Trong tuần này, ngày có biên độ nhiệt lớn nhất của Thành phố Hồ Chí Minh là thứ mấy?
	\item Chọn ngẫu nhiên một ngày trong tuần, tính xác suất của các biến cố sau:
	\begin{itemize}
		\item Biến cố $E$: ``Ngày được chọn có nhiệt độ cao nhất không quá $35^\circ\text{C}$''.
		\item Biến cố $F$: ``Ngày được chọn có biên độ nhiệt nhỏ hơn $12^\circ\text{C}$''.
	\end{itemize}
\end{enumerate}
\loigiai{
\begin{enumerate}[a)]
	\item Quan sát bảng số liệu biên độ nhiệt của các ngày trong tuần:
	\begin{itemize}
		\item Thứ Hai, Thứ Ba, Thứ Tư, Thứ Năm đều có biên độ nhiệt là $11^\circ\text{C}$.
		\item Thứ Sáu và Thứ Bảy đều có biên độ nhiệt lớn nhất là $13^\circ\text{C}$ ($35 - 22 = 13$ và $36 - 23 = 13$).
		\item Chủ Nhật có biên độ nhiệt là $7^\circ\text{C}$.
	\end{itemize}
	Vậy ngày có biên độ nhiệt lớn nhất trong tuần là Thứ Sáu và Thứ Bảy (đều đạt $13^\circ\text{C}$).
	\item Số các kết quả có thể xảy ra khi chọn ngẫu nhiên một ngày trong tuần là $n(\Omega) = 7$.
	\begin{itemize}
		\item Các ngày có nhiệt độ cao nhất không quá $35^\circ\text{C}$ (tức là $\le 35^\circ\text{C}$) gồm: Thứ Tư ($35^\circ\text{C}$), Thứ Sáu ($35^\circ\text{C}$) và Chủ Nhật ($34^\circ\text{C}$). Do đó có $3$ kết quả thuận lợi cho biến cố $E$.\\
		Xác suất của biến cố $E$ là: $P(E) = \dfrac{3}{7}$.
		\item Các ngày có biên độ nhiệt nhỏ hơn $12^\circ\text{C}$ gồm: Thứ Hai ($11^\circ\text{C}$), Thứ Ba ($11^\circ\text{C}$), Thứ Tư ($11^\circ\text{C}$), Thứ Năm ($11^\circ\text{C}$) và Chủ Nhật ($7^\circ\text{C}$). Do đó có $5$ kết quả thuận lợi cho biến cố $F$.\\
		Xác suất của biến cố $F$ là: $P(F) = \dfrac{5}{7}$.
	\end{itemize}
\end{enumerate}
}
\end{bt}

\begin{bt}%[9D6V21-3]
% nguon: Đề thi tham khảo Tuyển sinh 10 TP.HCM – Phụ lục Sở GD&ĐT TP.HCM – Bài 4
Một khu vườn hình chữ nhật có chiều dài và chiều rộng lần lượt là $70\text{ m}$ và $30\text{ m}$. Người ta dự tính mở rộng thêm khu vườn bằng cách cải tạo thêm $x\text{ (mét)}$ về mỗi phía ngoài của chiều dài và chiều rộng khu vườn.
\begin{enumerate}[a)]
	\item Viết biểu thức $S$ biểu diễn theo $x$ diện tích của khu vườn hình chữ nhật sau khi mở rộng.
	\item Biết rằng sau khi mở rộng thì diện tích của khu vườn lớn hơn diện tích ban đầu $1150\text{ m}^2$. Tìm giá trị của $x$ (làm tròn đến hàng phần mười của mét).
\end{enumerate}
\loigiai{
\begin{enumerate}[a)]
	\item Sau khi mở rộng thêm $x\text{ (m)}$ về cả hai phía ngoài của chiều dài và chiều rộng, khu vườn mới có:
	\begin{itemize}
		\item Chiều dài là: $70 + 2x\text{ (m)}$.
		\item Chiều rộng là: $30 + 2x\text{ (m)}$ (với $x > 0$).
	\end{itemize}
	Biểu thức biểu diễn diện tích khu vườn sau khi mở rộng là:
	$S = (70 + 2x)(30 + 2x) = 2100 + 140x + 60x + 4x^2 = 4x^2 + 200x + 2100\text{ (m}^2\text{)}$.
	\item Diện tích ban đầu của khu vườn là: $70\cdot 30 = 2100\text{ (m}^2\text{)}$.\\
	Theo đề bài, sau khi mở rộng thì diện tích khu vườn tăng thêm $1150\text{ m}^2$, nên ta có phương trình:
	$S - 2100 = 1150 \Leftrightarrow 4x^2 + 200x = 1150 \Leftrightarrow 4x^2 + 200x - 1150 = 0 \Leftrightarrow 2x^2 + 100x - 575 = 0$.\\
	Ta có $\Delta' = 50^2 - 2\cdot(-575) = 2500 + 1150 = 3650 > 0$.\\
	Vì $x > 0$ nên phương trình có nghiệm dương duy nhất là:
	$x = \dfrac{-50 + \sqrt{3650}}{2} \approx \dfrac{-50 + 60{,}415}{2} = \dfrac{10{,}415}{2} \approx 5{,}2\text{ (m)}$.\\
	Vậy giá trị của $x$ xấp xỉ bằng $5{,}2\text{ mét}$.
\end{enumerate}
}
\end{bt}

\begin{bt}%[9H0V31-3]
% nguon: Đề thi tham khảo Tuyển sinh 10 TP.HCM – Phụ lục Sở GD&ĐT TP.HCM – Bài 5
Một quả dưa hấu không hạt ruột đỏ dạng hình cầu có đường kính $25\text{ cm}$ và phần vỏ dày $2\text{ cm}$.
\begin{enumerate}[a)]
	\item Coi phần ruột màu đỏ cũng có dạng hình cầu có cùng tâm với quả dưa hấu. Tính thể tích phần ruột quả dưa hấu (kết quả làm tròn đến hàng phần trăm của $\text{cm}^3$).
	\item Người ta ép phần ruột màu đỏ của quả dưa hấu trên thì thể tích nước ép thu được bằng $80\%$ thể tích phần ruột. Nước ép dưa hấu sẽ được đựng trong các ly thuỷ tinh giống nhau, phần lòng trong dạng hình trụ có chiều cao $10\text{ cm}$ và đường kính đáy lòng trong là $5\text{ cm}$. Mỗi ly chỉ chứa được $70\%$ thể tích. Hỏi để đựng hết nước ép của quả dưa hấu nói trên thì cần ít nhất bao nhiêu cái ly?
\end{enumerate}
\loigiai{
\begin{enumerate}[a)]
	\item Bán kính của cả quả dưa hấu là: $R_1 = \dfrac{25}{2} = 12{,}5\text{ (cm)}$.\\
	Vì phần vỏ dưa dày $2\text{ cm}$ nên bán kính của phần ruột đỏ hình cầu là:
	$R = 12{,}5 - 2 = 10{,}5\text{ (cm)}$.\\
	Thể tích phần ruột quả dưa hấu là:
	$V_{\text{ruột}} = \dfrac{4}{3}\pi R^3 = \dfrac{4}{3}\pi (10{,}5)^3 = 1543{,}5\pi \approx 4849{,}05\text{ (cm}^3\text{)}$.
	\item Thể tích nước ép thu được từ quả dưa hấu là:
	$V_{\text{nước}} = 80\% \cdot V_{\text{ruột}} = 0{,}8 \cdot 1543{,}5\pi = 1234{,}8\pi\text{ (cm}^3\text{)}$.\\
	Bán kính đáy lòng trong của ly thuỷ tinh hình trụ là: $r = \dfrac{5}{2} = 2{,}5\text{ (cm)}$.\\
	Thể tích của mỗi chiếc ly là:
	$V_{\text{ly}} = \pi r^2 h = \pi \cdot (2{,}5)^2 \cdot 10 = 62{,}5\pi\text{ (cm}^3\text{)}$.\\
	Vì mỗi ly chỉ chứa được $70\%$ thể tích nên lượng nước ép đựng trong mỗi ly là:
	$V_{\text{chứa}} = 70\% \cdot V_{\text{ly}} = 0{,}7 \cdot 62{,}5\pi = 43{,}75\pi\text{ (cm}^3\text{)}$.\\
	Số lượng ly cần thiết để đựng hết lượng nước ép là:
	$N = \dfrac{V_{\text{nước}}}{V_{\text{chứa}}} = \dfrac{1234{,}8\pi}{43{,}75\pi} = \dfrac{1234{,}8}{43{,}75} = 28{,}224$.\\
	Vì số ly phải là số nguyên dương nên cần ít nhất $29$ chiếc ly để đựng hết lượng nước ép dưa hấu.
\end{enumerate}
}
\end{bt}

\begin{bt}%[9D1V3-4]
% nguon: Đề thi tham khảo Tuyển sinh 10 TP.HCM – Phụ lục Sở GD&ĐT TP.HCM – Bài 6
Thép không gỉ Ferritic là họ thép hợp kim có chứa từ $12\%$ đến $27\%$ crôm. Một nhà máy luyện thép hiện có sẵn một lượng hợp kim thép chứa $10\%$ crôm và một lượng hợp kim thép chứa $30\%$ crôm. Giả sử trong quá trình luyện thép các nguyên liệu không bị hao hụt.
\begin{enumerate}[a)]
	\item Tính khối lượng hợp kim thép mỗi loại từ hai loại thép trên dùng để luyện được $500\text{ tấn}$ thép chứa $16\%$ crôm.
	\item Nhà máy dự định luyện ra loại thép không gỉ Ferritic từ $100\text{ tấn}$ thép chứa $10\%$ crôm và $x\text{ tấn}$ thép chứa $30\%$ crôm. Hỏi $x$ nằm trong khoảng nào?
\end{enumerate}
\loigiai{
\begin{enumerate}[a)]
	\item Gọi khối lượng hợp kim thép chứa $10\%$ crôm và chứa $30\%$ crôm cần dùng lần lượt là $u$ và $v$ (tấn, điều kiện: $0 < u, v < 500$).\\
	Theo đề bài, tổng khối lượng thép luyện được là $500\text{ tấn}$, ta có:
	$u + v = 500$.\\
	Lượng crôm nguyên chất có trong sản phẩm là $500 \cdot 16\% = 80\text{ (tấn)}$, ta có phương trình:
	$10\% u + 30\% v = 80 \Leftrightarrow 0{,}1u + 0{,}3v = 80 \Leftrightarrow u + 3v = 800$.\\
	Ta có hệ phương trình:
	$\begin{cases} u + v = 500 \\ u + 3v = 800 \end{cases} \Leftrightarrow \begin{cases} 2v = 300 \\ u = 500 - v \end{cases} \Leftrightarrow \begin{cases} v = 150 \\ u = 350 \end{cases}$ (thoả mãn điều kiện).\\
	Vậy cần dùng $350\text{ tấn}$ thép chứa $10\%$ crôm và $150\text{ tấn}$ thép chứa $30\%$ crôm.
	\item Khi trộn $100\text{ tấn}$ thép chứa $10\%$ crôm với $x\text{ tấn}$ thép chứa $30\%$ crôm ($x > 0$):
	\begin{itemize}
		\item Tổng khối lượng thép thu được là: $100 + x\text{ (tấn)}$.
		\item Tổng lượng crôm nguyên chất thu được là: $100\cdot 10\% + x\cdot 30\% = 10 + 0{,}3x\text{ (tấn)}$.
	\end{itemize}
	Tỉ lệ phần trăm crôm trong hợp kim mới là: $\dfrac{10 + 0{,}3x}{100 + x}$.\\
	Để hợp kim thu được là thép không gỉ Ferritic (chứa từ $12\%$ đến $27\%$ crôm), ta có bất phương trình kép:
	$12\% \le \dfrac{10 + 0{,}3x}{100 + x} \le 27\% \Leftrightarrow 0{,}12 \le \dfrac{10 + 0{,}3x}{100 + x} \le 0{,}27$.\\
	Vì $x > 0$ nên $100 + x > 0$, nhân cả hai vế với $(100 + x)$:
	\begin{itemize}
		\item $0{,}12(100 + x) \le 10 + 0{,}3x \Leftrightarrow 12 + 0{,}12x \le 10 + 0{,}3x \Leftrightarrow 0{,}18x \ge 2 \Leftrightarrow x \ge \dfrac{2}{0{,}18} = \dfrac{100}{9} \approx 11{,}11\text{ (tấn)}$.
		\item $10 + 0{,}3x \le 0{,}27(100 + x) \Leftrightarrow 10 + 0{,}3x \le 27 + 0{,}27x \Leftrightarrow 0{,}03x \le 17 \Leftrightarrow x \le \dfrac{17}{0{,}03} = \dfrac{1700}{3} \approx 566{,}67\text{ (tấn)}$.
	\end{itemize}
	Vậy khối lượng $x$ cần dùng nằm trong khoảng $\dfrac{100}{9} \le x \le \dfrac{1700}{3}$ (tấn), tức là khoảng từ $11{,}11\text{ tấn}$ đến $566{,}67\text{ tấn}$.
\end{enumerate}
}
\end{bt}

\begin{bt}%[9H9C29-3]
% nguon: Đề thi tham khảo Tuyển sinh 10 TP.HCM – Phụ lục Sở GD&ĐT TP.HCM – Bài 7
Cho tam giác $ABC$ có ba góc nhọn ($AB < AC$). Đường tròn tâm $O$ đường kính $BC$ cắt hai cạnh $AB, AC$ lần lượt tại $E$ và $F$ ($E \neq B, F \neq C$). Các đoạn thẳng $BF$ và $CE$ cắt nhau tại $H$, tia $AH$ cắt $BC$ tại $K$.
\begin{enumerate}[a)]
	\item Chứng minh $\widehat{BEC} = \widehat{BFC} = 90^\circ$, từ đó suy ra tứ giác $AEHF$ nội tiếp.
	\item Gọi $D$ là giao điểm của $AH$ và đường tròn $(O)$ ($D$ nằm giữa $A$ và $H$). Chứng minh $BD^2 = BK \cdot BC$ và $\widehat{BDH} = \widehat{BFD}$.
	\item Trong trường hợp $\widehat{BAC} = 60^\circ$ và $BC = 6\text{ cm}$, tính độ dài đoạn thẳng $EF$ và bán kính đường tròn ngoại tiếp tam giác $AEF$.
\end{enumerate}
\loigiai{
\begin{enumerate}[a)]
	\item Điểm $E, F$ thuộc đường tròn $(O)$ đường kính $BC$ nên $\widehat{BEC} = 90^\circ$ và $\widehat{BFC} = 90^\circ$ (các góc nội tiếp chắn nửa đường tròn).\\
	Suy ra $CE \perp AB$ và $BF \perp AC$. Do đó $CE$ và $BF$ là hai đường cao của tam giác $ABC$.\\
	Vì $BF$ cắt $CE$ tại $H$ nên $H$ là trực tâm của tam giác $ABC$, suy ra $AH \perp BC$ tại $K$, hay $\widehat{AKC} = 90^\circ$.\\
	Xét tứ giác $AEHF$, ta có $\widehat{AEH} = 90^\circ$ (do $CE \perp AB$) và $\widehat{AFH} = 90^\circ$ (do $BF \perp AC$).\\
	Suy ra $\widehat{AEH} + \widehat{AFH} = 90^\circ + 90^\circ = 180^\circ$.\\
	Vậy tứ giác $AEHF$ nội tiếp đường tròn đường kính $AH$.
	\item Điểm $D$ thuộc đường tròn $(O)$ đường kính $BC$ nên tam giác $BDC$ vuông tại $D$, tức là $\widehat{BDC} = 90^\circ$.\\
	Vì $AH \perp BC$ tại $K$ nên $DK \perp BC$.\\
	Xét tam giác $BDC$ vuông tại $D$ có đường cao $DK$. Áp dụng hệ thức lượng trong tam giác vuông, ta có:
	$BD^2 = BK \cdot BC$.\\
	Xét tam giác $BDK$ và tam giác $BCD$, ta có góc $\widehat{CBD}$ chung và $\dfrac{BD}{BK} = \dfrac{BC}{BD}$ (do $BD^2 = BK \cdot BC$), suy ra $\Delta BDK \backsim \Delta BCD$ (c.g.c).\\
	Do đó $\widehat{BDK} = \widehat{BCD}$, hay $\widehat{BDH} = \widehat{BCD}$.\\
	Mặt khác, trong đường tròn $(O)$, góc $\widehat{BFD}$ và góc $\widehat{BCD}$ là hai góc nội tiếp cùng chắn cung $BD$, nên $\widehat{BFD} = \widehat{BCD}$.\\
	Từ đó suy ra $\widehat{BDH} = \widehat{BFD}$.
	\item Xét tam giác $AEF$ và tam giác $ACB$:
	Trong tam giác vuông $ABF$ tại $F$, ta có: $\cos A = \dfrac{AF}{AB}$.\\
	Trong tam giác vuông $ACE$ tại $E$, ta có: $\cos A = \dfrac{AE}{AC}$.\\
	Suy ra $\dfrac{AE}{AC} = \dfrac{AF}{AB} = \cos 60^\circ = \dfrac{1}{2}$.\\
	Xét tam giác $AEF$ và tam giác $ACB$, có góc $\widehat{BAC}$ chung và $\dfrac{AE}{AC} = \dfrac{AF}{AB}$, do đó $\Delta AEF \backsim \Delta ACB$ (c.g.c) theo tỉ số đồng dạng $k = \cos A = \dfrac{1}{2}$.\\
	Suy ra $\dfrac{EF}{BC} = k = \dfrac{1}{2} \Rightarrow EF = \dfrac{1}{2}BC = \dfrac{1}{2}\cdot 6 = 3\text{ (cm)}$.\\
	Vì tứ giác $AEHF$ nội tiếp đường tròn đường kính $AH$ nên đường tròn ngoại tiếp tam giác $AEF$ chính là đường tròn đường kính $AH$.\\
	Do đó bán kính đường tròn ngoại tiếp tam giác $AEF$ là $R_{AEF} = \dfrac{AH}{2}$.\\
	Gọi $R_0$ là bán kính đường tròn ngoại tiếp tam giác $ABC$. Áp dụng định lí sin cho tam giác $ABC$:
	$\dfrac{BC}{\sin A} = 2R_0 \Rightarrow 2R_0 = \dfrac{6}{\sin 60^\circ} = \dfrac{6}{\frac{\sqrt{3}}{2}} = 4\sqrt{3} \Rightarrow R_0 = 2\sqrt{3}\text{ (cm)}$.\\
	Mặt khác, trong một tam giác nhọn, khoảng cách từ trực tâm $H$ đến đỉnh $A$ bằng $AH = 2R_0 \cos A$.\\
	Do đó: $AH = 2 \cdot (2\sqrt{3}) \cdot \cos 60^\circ = 4\sqrt{3} \cdot \dfrac{1}{2} = 2\sqrt{3}\text{ (cm)}$.\\
	Suy ra bán kính đường tròn ngoại tiếp tam giác $AEF$ là:
	$R_{AEF} = \dfrac{AH}{2} = \dfrac{2\sqrt{3}}{2} = \sqrt{3}\text{ (cm)}$.\\
	Vậy $EF = 3\text{ cm}$ và bán kính đường tròn ngoại tiếp tam giác $AEF$ bằng $\sqrt{3}\text{ cm}$.
\end{enumerate}
}
\end{bt}
